%% ODER: format ==         = "\mathrel{==}"
%% ODER: format /=         = "\neq "
%
%
\makeatletter
\@ifundefined{lhs2tex.lhs2tex.sty.read}%
  {\@namedef{lhs2tex.lhs2tex.sty.read}{}%
   \newcommand\SkipToFmtEnd{}%
   \newcommand\EndFmtInput{}%
   \long\def\SkipToFmtEnd#1\EndFmtInput{}%
  }\SkipToFmtEnd

\newcommand\ReadOnlyOnce[1]{\@ifundefined{#1}{\@namedef{#1}{}}\SkipToFmtEnd}
\usepackage{amstext}
\usepackage{amssymb}
\usepackage{stmaryrd}
\DeclareFontFamily{OT1}{cmtex}{}
\DeclareFontShape{OT1}{cmtex}{m}{n}
  {<5><6><7><8>cmtex8
   <9>cmtex9
   <10><10.95><12><14.4><17.28><20.74><24.88>cmtex10}{}
\DeclareFontShape{OT1}{cmtex}{m}{it}
  {<-> ssub * cmtt/m/it}{}
\newcommand{\texfamily}{\fontfamily{cmtex}\selectfont}
\DeclareFontShape{OT1}{cmtt}{bx}{n}
  {<5><6><7><8>cmtt8
   <9>cmbtt9
   <10><10.95><12><14.4><17.28><20.74><24.88>cmbtt10}{}
\DeclareFontShape{OT1}{cmtex}{bx}{n}
  {<-> ssub * cmtt/bx/n}{}
\newcommand{\tex}[1]{\text{\texfamily#1}}	% NEU

\newcommand{\Sp}{\hskip.33334em\relax}


\newcommand{\Conid}[1]{\mathit{#1}}
\newcommand{\Varid}[1]{\mathit{#1}}
\newcommand{\anonymous}{\kern0.06em \vbox{\hrule\@width.5em}}
\newcommand{\plus}{\mathbin{+\!\!\!+}}
\newcommand{\bind}{\mathbin{>\!\!\!>\mkern-6.7mu=}}
\newcommand{\rbind}{\mathbin{=\mkern-6.7mu<\!\!\!<}}% suggested by Neil Mitchell
\newcommand{\sequ}{\mathbin{>\!\!\!>}}
\renewcommand{\leq}{\leqslant}
\renewcommand{\geq}{\geqslant}
\usepackage{polytable}

%mathindent has to be defined
\@ifundefined{mathindent}%
  {\newdimen\mathindent\mathindent\leftmargini}%
  {}%

\def\resethooks{%
  \global\let\SaveRestoreHook\empty
  \global\let\ColumnHook\empty}
\newcommand*{\savecolumns}[1][default]%
  {\g@addto@macro\SaveRestoreHook{\savecolumns[#1]}}
\newcommand*{\restorecolumns}[1][default]%
  {\g@addto@macro\SaveRestoreHook{\restorecolumns[#1]}}
\newcommand*{\aligncolumn}[2]%
  {\g@addto@macro\ColumnHook{\column{#1}{#2}}}

\resethooks

\newcommand{\onelinecommentchars}{\quad-{}- }
\newcommand{\commentbeginchars}{\enskip\{-}
\newcommand{\commentendchars}{-\}\enskip}

\newcommand{\visiblecomments}{%
  \let\onelinecomment=\onelinecommentchars
  \let\commentbegin=\commentbeginchars
  \let\commentend=\commentendchars}

\newcommand{\invisiblecomments}{%
  \let\onelinecomment=\empty
  \let\commentbegin=\empty
  \let\commentend=\empty}

\visiblecomments

\newlength{\blanklineskip}
\setlength{\blanklineskip}{0.66084ex}

\newcommand{\hsindent}[1]{\quad}% default is fixed indentation
\let\hspre\empty
\let\hspost\empty
\newcommand{\NB}{\textbf{NB}}
\newcommand{\Todo}[1]{$\langle$\textbf{To do:}~#1$\rangle$}

\EndFmtInput
\makeatother
%
%
%
%
%
%
% This package provides two environments suitable to take the place
% of hscode, called "plainhscode" and "arrayhscode". 
%
% The plain environment surrounds each code block by vertical space,
% and it uses \abovedisplayskip and \belowdisplayskip to get spacing
% similar to formulas. Note that if these dimensions are changed,
% the spacing around displayed math formulas changes as well.
% All code is indented using \leftskip.
%
% Changed 19.08.2004 to reflect changes in colorcode. Should work with
% CodeGroup.sty.
%
\ReadOnlyOnce{polycode.fmt}%
\makeatletter

\newcommand{\hsnewpar}[1]%
  {{\parskip=0pt\parindent=0pt\par\vskip #1\noindent}}

% can be used, for instance, to redefine the code size, by setting the
% command to \small or something alike
\newcommand{\hscodestyle}{}

% The command \sethscode can be used to switch the code formatting
% behaviour by mapping the hscode environment in the subst directive
% to a new LaTeX environment.

\newcommand{\sethscode}[1]%
  {\expandafter\let\expandafter\hscode\csname #1\endcsname
   \expandafter\let\expandafter\endhscode\csname end#1\endcsname}

% "compatibility" mode restores the non-polycode.fmt layout.

\newenvironment{compathscode}%
  {\par\noindent
   \advance\leftskip\mathindent
   \hscodestyle
   \let\\=\@normalcr
   \let\hspre\(\let\hspost\)%
   \pboxed}%
  {\endpboxed\)%
   \par\noindent
   \ignorespacesafterend}

\newcommand{\compaths}{\sethscode{compathscode}}

% "plain" mode is the proposed default.
% It should now work with \centering.
% This required some changes. The old version
% is still available for reference as oldplainhscode.

\newenvironment{plainhscode}%
  {\hsnewpar\abovedisplayskip
   \advance\leftskip\mathindent
   \hscodestyle
   \let\hspre\(\let\hspost\)%
   \pboxed}%
  {\endpboxed%
   \hsnewpar\belowdisplayskip
   \ignorespacesafterend}

\newenvironment{oldplainhscode}%
  {\hsnewpar\abovedisplayskip
   \advance\leftskip\mathindent
   \hscodestyle
   \let\\=\@normalcr
   \(\pboxed}%
  {\endpboxed\)%
   \hsnewpar\belowdisplayskip
   \ignorespacesafterend}

% Here, we make plainhscode the default environment.

\newcommand{\plainhs}{\sethscode{plainhscode}}
\newcommand{\oldplainhs}{\sethscode{oldplainhscode}}
\plainhs

% The arrayhscode is like plain, but makes use of polytable's
% parray environment which disallows page breaks in code blocks.

\newenvironment{arrayhscode}%
  {\hsnewpar\abovedisplayskip
   \advance\leftskip\mathindent
   \hscodestyle
   \let\\=\@normalcr
   \(\parray}%
  {\endparray\)%
   \hsnewpar\belowdisplayskip
   \ignorespacesafterend}

\newcommand{\arrayhs}{\sethscode{arrayhscode}}

% The mathhscode environment also makes use of polytable's parray 
% environment. It is supposed to be used only inside math mode 
% (I used it to typeset the type rules in my thesis).

\newenvironment{mathhscode}%
  {\parray}{\endparray}

\newcommand{\mathhs}{\sethscode{mathhscode}}

% texths is similar to mathhs, but works in text mode.

\newenvironment{texthscode}%
  {\(\parray}{\endparray\)}

\newcommand{\texths}{\sethscode{texthscode}}

% The framed environment places code in a framed box.

\def\codeframewidth{\arrayrulewidth}
\RequirePackage{calc}

\newenvironment{framedhscode}%
  {\parskip=\abovedisplayskip\par\noindent
   \hscodestyle
   \arrayrulewidth=\codeframewidth
   \tabular{@{}|p{\linewidth-2\arraycolsep-2\arrayrulewidth-2pt}|@{}}%
   \hline\framedhslinecorrect\\{-1.5ex}%
   \let\endoflinesave=\\
   \let\\=\@normalcr
   \(\pboxed}%
  {\endpboxed\)%
   \framedhslinecorrect\endoflinesave{.5ex}\hline
   \endtabular
   \parskip=\belowdisplayskip\par\noindent
   \ignorespacesafterend}

\newcommand{\framedhslinecorrect}[2]%
  {#1[#2]}

\newcommand{\framedhs}{\sethscode{framedhscode}}

% The inlinehscode environment is an experimental environment
% that can be used to typeset displayed code inline.

\newenvironment{inlinehscode}%
  {\(\def\column##1##2{}%
   \let\>\undefined\let\<\undefined\let\\\undefined
   \newcommand\>[1][]{}\newcommand\<[1][]{}\newcommand\\[1][]{}%
   \def\fromto##1##2##3{##3}%
   \def\nextline{}}{\) }%

\newcommand{\inlinehs}{\sethscode{inlinehscode}}

% The joincode environment is a separate environment that
% can be used to surround and thereby connect multiple code
% blocks.

\newenvironment{joincode}%
  {\let\orighscode=\hscode
   \let\origendhscode=\endhscode
   \def\endhscode{\def\hscode{\endgroup\def\@currenvir{hscode}\\}\begingroup}
   %\let\SaveRestoreHook=\empty
   %\let\ColumnHook=\empty
   %\let\resethooks=\empty
   \orighscode\def\hscode{\endgroup\def\@currenvir{hscode}}}%
  {\origendhscode
   \global\let\hscode=\orighscode
   \global\let\endhscode=\origendhscode}%

\makeatother
\EndFmtInput
%
%
%
% First, let's redefine the forall, and the dot.
%
%
% This is made in such a way that after a forall, the next
% dot will be printed as a period, otherwise the formatting
% of `comp_` is used. By redefining `comp_`, as suitable
% composition operator can be chosen. Similarly, period_
% is used for the period.
%
\ReadOnlyOnce{forall.fmt}%
\makeatletter

% The HaskellResetHook is a list to which things can
% be added that reset the Haskell state to the beginning.
% This is to recover from states where the hacked intelligence
% is not sufficient.

\let\HaskellResetHook\empty
\newcommand*{\AtHaskellReset}[1]{%
  \g@addto@macro\HaskellResetHook{#1}}
\newcommand*{\HaskellReset}{\HaskellResetHook}

\global\let\hsforallread\empty

\newcommand\hsforall{\global\let\hsdot=\hsperiodonce}
\newcommand*\hsperiodonce[2]{#2\global\let\hsdot=\hscompose}
\newcommand*\hscompose[2]{#1}

\AtHaskellReset{\global\let\hsdot=\hscompose}

% In the beginning, we should reset Haskell once.
\HaskellReset

\makeatother
\EndFmtInput
% ----------------------
%% Stuff for TimCore


%% 'g' inferences




%% --------- Effects ------------

%% Intersection of effects




%% ----------------------


%% --------------------------
%%        Desugaring rules
%% --------------------------



% ----------------------------------------------------
% New desugaring, just does the wrapping part:
% ----------------------------------------------------
%   Input for wrapping


%%format Hide (e) = "{\mathcal H}\llbracket{}" e "\rrbracket{}"

%% 'si' is short for "solve or infer" or "unification variable or not" resp

%% 'md' is short for "mode" or "unification variable or not" resp

%% flipsi is a no-op for now


%% EXX is just EX, but with no argument




%% ----------- Applications --------------
%% %format applyx a b = a "@" b
%% %format applyx a b = a "\," b

%% Function application with no arguments

%% Function application with 1 argument

%% Function application with 1 argument, no parens

%% Function application with 1 argument, without parenthesis

%% Function application with many arguments

%% ----------- End of Applications --------------





%% Sequences e1; ...; en; v

%% xs etc stole syntax used in examples.

%% Old version: | for BNF, [] for choice
%%%format choice = "[ \! ]"
%%%format `choice` = "\mathop{[ \! ]}"
%% 1mm is about 0.3em with the current font

%% New version: tall | for BNF, fat | for choice
%%format vbar = "\mathop{\bigm|}"
%%format choice = "\boldsymbol{|}"
%%format `choice` = "\mathop{\boldsymbol{|}}"




%% %format ok = "\textbf{ok}"

%% %format `eqsemi` = "\hspace{-0.5mm}\fatsemi{}"
%% %format eqsemiKW = "\fatsemi"

%% Arrays and tuples

%% %format vmap (x) = "\llangle" x "\rrangle"





%%  OLD SYNTAX   %format split (x) (y) (z) = "\textbf{split}(" x ")\{" y "," z "\}"
%% %format defKW = "\textbf{def}"
% only used in comments, but still prettier this way:
%%%%format map = "\textbf{map}"

%% -------Functions ---------------


%% functionOC has open/closed subscript

%   \_q<fx>i.(e1)(y:=f[x]){e2}

%% -------End of functions ---------------





%% Unification
%%%format == = "\!=\!"


% format := = "\mapsto"
% %format squiggt = "\leadsto"
% %format squiggt = "\!-\!\!>\!\!\!"






























%% ----------- DPS Verse -------------

%% Translation from Mini Verse to DPS Verse


%% ----------- End of DPS Verse -------------






%% We combined S and C into one


%% ----------- Metatheory --------------


%% ----------------------------
%%         Verification
%% ----------------------------

%% assert = succeed, abbreviated to suc



%% Effects checking, like succeeds{e}





%% <fx> brackets

% -------------------------------------------------------------
%               The ENV type and its operations
% -------------------------------------------------------------

\begin{figure}
$$
\begin{array}{l}

\textbf{Domains}\\
\begin{array}{rl@{\;\;}c@{\;\;}ll}
  v \in & \Value & = & \Int(\Z) + \Fun(\F) \\\
  & \F & = & \Vm{\semfcn}  & \text{Any invariants? e.g. disjoint domains} \\
  \rho \in & \Env & = & \Ident \rightarrow \Value  & \text{Total: gives a value to every identifier} \\
  \EV \in &\ENV & = & \powerset{ \Env } \\
  \{ \fpr{x_1}{y_1}, \mydots, \fpr{x_n}{y_n} \} \in & \semfcn & = & \powerset{ \Val \times \Val }
  & \text{Partial functions} \\
\end{array} \\
\\
\textbf{The constraint language}\\
\begin{array}{l@{\;\;}rc@{\;\;}ll}
  \text{Simple values}    & a & ::= & \ensuremath{\Varid{v}} \vb \ensuremath{\Varid{x}} \vb op[a_1,\mydots,a_n] \vb \ensuremath{\langle\Varid{a}_{\mathrm{0}},\mycdots,\Varid{a}_{\Varid{n}}\rangle} \\
% \vb |truth a|\\
  \text{Constraints} &  \EV & ::= & \conBrack{relop[a_1,\mydots,a_n]} \\
   &  & \vb & \conBrack{ \fpr{x}{y} \in h } \vb \conBrack{ x \not\in \dom(h) }  & \text{where}\; h \in \semfcn \\
%   &  & \vb & \conBrack{ \conFun{a_{arg}}{a_{res}}{a_{fun}}{i}{n}}  \\
  &  & \vb & \conEmpty \vb \conUniv \\
  &  & \vb & \EV_1 \cup \EV_2 \vb \EV_1 \cap \EV_2 \\
  &  & \vb & \conCompl(\EV)  & \text{Complement} \\
  &  & \vb & \HIDE{\EV}{xs} & \text{Hiding} \\
  \text{Relational ops} &  relop & ::= & = \vb < \vb \leq \vb \geq \vb > \\
  \text{Value ops} & op & ::= & + \vb - \vb * \vb / \\
  & & \vb & \funcat \vb \tupcat  & \text{Function and tuple concatenation resp.} \\
  & & \vb & \fununion & \text{Function/tuple union} \\
\end{array}\\
\\[-1mm]
\textbf{Interpretation of simple values and constraints}\\
\begin{array}{r@{\;\;}c@{\;\;}ll}
  \sembrack{\EV} & : & \ENV \\
  \sembrack{\conEmpty} & = & \{\} \\
  \sembrack{\conUniv} & = & \Env   & \text{Set of all environments} \\
  \sembrack{\conCompl(\EV)} & = & \complE( \sembrack{\EV} )\\
  \sembrack{\HIDE{\EV}{xs}} & = & \{ \rho \in \Env \vb \exists \rho' \in \sembrack{\EV} .\, \forall x \not\in xs.\, \rho(x) =  \rho'(x) \} \\
  \sembrack{\EV_1 \cup \EV_2} & = & \sembrack{\EV_1} \cup \sembrack{\EV_2} \\
  \sembrack{\EV_1 \cap \EV_2} & = & \sembrack{\EV_1} \cap \sembrack{\EV_2} \\
  \sembrack{\conBrack{a_1=a_2}} & = & \{ \rho \in \Env \vb \semapp{A}{a_1}\rho = \semapp{A}{a_2}\rho \} \\
  \sembrack{\conBrack{a_1>a_2}} & = & \{ \rho \in \Env \vb \semapp{A}{a_1}\rho > \semapp{A}{a_2}\rho \}
      \hspace{4mm} \text{...and similarly other $relops$} \\
  \sembrack{ \conBrack{\fpr{x}{y} \in h} } & = & \{ \rho \in \Env \vb \fpr{\rho(x)}{\rho(y)} \in h \} \\
  \sembrack{ \conBrack{ x \not\in \dom(h) } } & = & \{ \rho \in \Env \vb \rho(x) \not\in \dom(h) \}
\end{array} \\
\\
\begin{array}{r@{\;\;}c@{\;\;}ll}
\semapp{A}{\ensuremath{\Varid{a}}} & : & \parfun{\Env}{\Val} \\
\semapp{A}{\ensuremath{\Varid{x}}}\rho & = & \rho(x) \\
\semapp{A}{\ensuremath{\Varid{v}}}\rho & = & v \\
%\semapp{A}{|o|}\rho & = & \semap{O}{o} \\
\semapp{A}{\ensuremath{\Varid{a}_{\mathrm{1}}\mathbin{+}\Varid{a}_{\mathrm{2}}}}\rho & = & \Int(r_1 + r_2)
  & \text{if $\Int(r_1)= \semapp{A}{a_1}\rho$ and
     $\Int(r_2) = \semapp{A}{a_2}\rho$ } \\

% \semapp{A}{|truth a|}\rho & = &  FUN [ \fpr{v}{v} ] \hspace{5mm} \text{where $v = \semapp{A}{a}\rho$} \\
\semapp{A}{\ensuremath{\Varid{a}_{\mathrm{1}}} \fununion \ensuremath{\Varid{a}_{\mathrm{2}}}}\rho & = & \Fun(r_1 \pomunion r_2)
    &\text{if $\Fun(r_1)= \semapp{A}{a_1}\rho$, $\Fun(r_2)= \semapp{A}{a_2}\rho$,
                      $\pomdom(r_1)\cap \pomdom(r_2) = \emptyset$} \\
\semapp{A}{\ensuremath{\Varid{a}_{\mathrm{1}}} \funcat \ensuremath{\Varid{a}_{\mathrm{2}}}}\rho & = & \Fun(r_1 \pomcat r_2)
    & \text{if $\Fun(r_1)= \semapp{A}{a_1}\rho$, $\Fun(r_2)= \semapp{A}{a_2}\rho$,
                      $\pomdom(r_1)\cap \pomdom(r_2) = \emptyset$} \\
\semapp{A}{\ensuremath{\Varid{a}_{\mathrm{1}}} \tupcat \ensuremath{\Varid{a}_{\mathrm{2}}}}\rho & = & \Fun(\pomtupcat(r_1,r_2))
    & \text{if $\Fun(r_1)= \semapp{A}{a_1}\rho$, $\Fun(r_2)= \semapp{A}{a_2}\rho$,
                      $\pomdom(r_1)\cap \pomdom(r_2) = \emptyset$} \\
\semapp{A}{\ensuremath{\langle\rangle}}\rho & = & \Fun(\pomunit{ \pomempty } ) \\
\semapp{A}{\ensuremath{\langle\Varid{a}_{\mathrm{0}},\mycdots,\Varid{a}_{\Varid{n}}\rangle}}\rho
  & = & \multicolumn{2}{@{}l}{\Fun ( \; \pomunit{ \{ \fpr{0}{\semapp{A}{a_0}\rho} \}} \pomcat \ldots \pomcat
            \pomunit{ \{ \fpr{n}{\semapp{A}{a_n}\rho} \}} \; )} \\
% \semap{O}{|o|} & : & \Val \\
% \semap{O}{|<=|} & = & \mul\Fun[\,\{ \fpr{\tupE(x,y)}{x} \vb x \in \Z, y \in \Z, x \leq y \}\,] \\[1mm]
\end{array}\\
\\[-1mm]
\textbf{Semantic functions} \\

\begin{array}{lr@{\;\;}c@{\;\;}l}
\text{Binders of block} & \semap{I}{\ensuremath{\Varid{t}}} & : & \powerset{\Ident} \\
&  \semap{I}{\ensuremath{\Varid{t}}} & = & \text{the outer $\exists$-bound variables in \ensuremath{\Varid{t}}} \\[1mm]

% \text{Tuple values} &  \tupE & : & [\Val] \rightarrow \Val \\
% &  \tupE [v_0,\mydots,v_n] & = & \Fun \pomcompL \{ \fpr{i}{v_i} \} \vb i |<-| 0 \mycdots n \pomcompR
%   \hspace{1cm} \text{So $\tupE [] = \Fun(\pomempty)$}
\end{array} \\

\end{array}
$$
\caption{Semantic domains and associated functions}\label{fig:densem-notation}
\end{figure}


% -------------------------------------------------------------
%               The Verse Monad
% -------------------------------------------------------------

\begin{figure}
$$
\begin{array}{l}
\text{The Verse monad $\Vm{\cdot}$ comes equipped with the following operations} \\
  \begin{array}{rcll}
    \vmempty & :: & \Vm{a} \\
    \vmreturn & :: & a \rightarrow \Vm{a} \\
%    (\vmbind) & :: & \Vm{a} \rightarrow (a \rightarrow \Vm{b}) \rightarrow \Vm{b} \\
    \vminject & :: & \powerset{a} \rightarrow \Vm{a} \\
    \vmcat    & :: & \Vm{a} \rightarrow \Vm{a} \rightarrow \Vm{a} & \text{Ordered choice} \\
    \vmseq    & :: & \Vm{a} \rightarrow \Vm{a} \rightarrow \Vm{a} & \text{Inner-set intersection}\\
    \vminnerunion  & :: & \Vm{a} \rightarrow \Vm{a} \rightarrow \Vm{a} & \text{Inner-set union} \\
    \vmflatten & :: & \Vm{\powerset{a}} \rightarrow \Vm{a} \\
    \vmunion  & :: & \Vm{a} \rightarrow \Vm{a} \rightarrow \Vm{a} & \text{Unordered choice}\\
    \vmunions & :: & \powerset{\Vm{a}} \rightarrow \Vm{a}  & \text{Infinite unordered choice} \\
    \vmmap    & :: & (a \rightarrow b) \rightarrow \Vm{a} \rightarrow \Vm{b} \\
    \vmfold   & :: & \multicolumn{2}{l}{(\powerset{a} \rightarrow \Vm{b} \rightarrow \Vm{b})
                          \rightarrow \Vm{b} \rightarrow \Vm{a} \rightarrow \Vm{b}} \\
    \vmpi & : & \Vm{a} \rightarrow (a \rightarrow \powerset{b}) \rightarrow \powerset{\Vm{a,b}} \\
    \vmcollapse & :: & \Vm{a} \rightarrow \powerset{a} \\
    \vmnot      & :: & \Vm{a} \rightarrow \Vm{a}
  \end{array}
  \\ \\
  \textbf{Desirable laws} \\
  \begin{array}{rclcl}
%     \vmempty \vmbind k & = & \vmempty \\
     \vmempty \vmseq   s & = & \vmempty & = & s \vmseq \vmempty \\
     \vmempty \vmcat   s & = & s        & = & s \vmcat \vmempty \\
     \vmempty \vmunion s & = & s        & = & s \vmunion \vmempty \\
     \vmcollapse(\vminject(s)) & = & s \\
  \end{array}
\\ \\
%   \textbf{Monad comprehensions} \\
%  \text{The Verse monad enjoys the usual monad-comprehension notation:} \\
%  \begin{array}{rcl}
%    \vmcompL e \vmcompR & = & \vmreturn(e) \\
%    \vmcompL e_1 \vb e_2, \; Q \vmcompR
%       & = & |if| \; e_2 \; |then| \; \vmcompL e_1 \vb Q \vmcompR \; |else|\; \vmempty \\
%    \vmcompL e_1 \vb x \leftarrow e_2, \; Q \vmcompR & = & e_2 \vmbind \lambda x. \vmcompL e_1 \vb Q \vmcompR \\
%    \vmcompL e_1 \vb x \in e_2, \; Q \vmcompR & = & \vminject(e_2) \vmbind \lambda x. \vmcompL e_1 \vb Q \vmcompR \\
%  \end{array}
% \\ \\
  \textbf{Derived functions} \\
  \begin{array}{rcl}
    \vmflatmap & :: & (a \rightarrow \powerset{b}) \rightarrow \Vm{a} \rightarrow \Vm{b} \\
    \vmflatmap(f,s) & = & \vmflatten(\; \vmmap( f, s ) \; ) \\[1mm]
    (\ENVhide) & :: & \powerset{\Env} \rightarrow \powerset{\Ident} \rightarrow \powerset{\Env} \\
    \rho{}s \ENVhide vs & = & \{ \rho' \in \Env \vb \exists \rho \in \rho{}s, \forall x \not\in vs \Rightarrow \rho(x) = \rho'(x) \} \\[1mm]
    (\vmhide) & :: & \Vm{\Env} \rightarrow \powerset{\Ident} \rightarrow \Vm{\Env} \\
    s \vmhide vs & = & \vmflatmap (\; \lambda \rho. \{ \rho \} \ENVhide vs, \;s\; ) \\[1mm]
    \vmone & :: & \Vm{\Env} \rightarrow \powerset{\Ident} \rightarrow \Vm{\Env} \\
    \vmone( s, vs ) & = & \vmfold( op, \vminject\{\}, s ) \\
    & \multicolumn{2}{l}{\text{where}} \\
    && \begin{array}[t]{@{}l}
         op :: \powerset{\Env} \rightarrow \Vm{\Env} \rightarrow \Vm{\Env} \\
         op( \rho{}s, rest ) = \vminject( \rho{}s ) \vminnerunion (\vminject( \conCompl( \rho{}s \ENVhide vs )) \vmseq rest )
       \end{array} \\
\\
    \vmconcat & :: & [\Vm{a}] \rightarrow \Vm{a} \\
    \vmconcat(\, [] \, ) & = & \vmempty \\
    \vmconcat(\, s:ss \, ) & = & s \vmcat \vmconcat(ss) \\
    \end{array}

\end{array}
$$
\caption{The Verse monad}
\end{figure}

% ------------------------------------------------------------------------------------
%
%            The semantics itself
%
% ------------------------------------------------------------------------------------

\begin{figure}
$$
\begin{array}{l}
\begin{array}{@{}lr@{\;\;}c@{\;\;}ll}
  & \semess{E}{u}{v}{\ensuremath{\Varid{t}}} & : & \VmENV  \\

% Wildcards and atoms
\rulename{swild}  &  \semess{E}{u}{v}{\ensuremath{\anonymous }} & = & \vminject \conBrack{u =v} \\
\rulename{svar}   &  \semess{E}{u}{v}{\ensuremath{\Varid{x}}} & = & \vminject \conBrack{u=v=x} \\
\rulename{sconst} &  \semess{E}{u}{v}{\ensuremath{\Varid{k}}} & = & \vminject \conBrack{u=v=k}\\
\rulename{sprim}  &  \semess{E}{u}{v}{\ensuremath{\Varid{o}}} & = & \vminject \conBrack{u=v=o} \\

% Capture input and output
\rulename{sbind} &  \semess{E}{u}{v}{\ensuremath{\Varid{x}\mathrel{\coloneqq}\Varid{t}}} & = &
   \vminject\conBrack{x=v} \pomseq \semess{E}{u}{v}{t} \\
\rulename{ssquig} &  \semess{E}{u}{v}{\ensuremath{\Varid{x}\!:\shortrightarrow\!\Varid{t}}} & = &
   \vminject\conBrack{x=u} \pomseq \semess{E}{u}{v}{t} \\

% Tuples
\multicolumn{2}{@{}l}{\rulename{stup} \; \hfill \semess{E}{u}{v}{\ensuremath{\textbf{array}\{\Varid{t}_{\mathrm{1}},\mycdots,\Varid{t}_{\Varid{n}}\}}}}
  & = & (\vminject \conBrack{ u=\ensuremath{\langle\Varid{u}_{\mathrm{1}},\mycdots,\Varid{u}_{\Varid{n}}\rangle}, v=\ensuremath{\langle\Varid{v}_{\mathrm{1}},\mycdots,\Varid{v}_{\Varid{n}}\rangle}} \\
  &&&   \; \vmseq \semess{E}{u_1}{v_1}{t_1}
         \vmseq \mydots
         \vmseq \semess{E}{u_n}{v_n}{t_n}) \\
&&&      \vmhide { \{u_1,\mydots u_n,v_1 \mydots v_n\} }
         \hfill u_i,v_i\not\in \freevars{t_1 \mydots t_n, u,v} \\

% Choice and failure
\rulename{sfail} &  \semess{E}{u}{v}{\ensuremath{\textbf{fail}}} & = & \vmempty \\
\rulename{schoice} &  \semess{E}{u}{v}{\ensuremath{\Varid{t}_{\mathrm{1}}\hspace{0.19em}\mathop{\rule[-0.1em]{0.15em}{0.75em}}\hspace{0.2em}\Varid{t}_{\mathrm{2}}}}
  & = & \semess{B}{u}{v}{t_1} \vmcat \semess{B}{u}{v}{t_2} \\
\rulename{suchoice} &  \semess{E}{u}{v}{\ensuremath{\Varid{t}_{\mathrm{1}}|\!|\!|\Varid{t}_{\mathrm{2}}}}
  & = & \semess{B}{u}{v}{t_1} \vmunion \semess{B}{u}{v}{t_2} \\

% Unification
\rulename{sunify} &  \semess{E}{u}{v}{\ensuremath{\Varid{t}_{\mathrm{1}}\hspace{-0.14em}=\hspace{-0.14em}\Varid{t}_{\mathrm{2}}}} & = &   \semess{E}{u}{v} {t_1} \vmseq \semess{E}{u}{v}{t_2} \\

% Sequencing
\rulename{sseq} &  \semess{E}{u}{v}{\ensuremath{\Varid{t}_{\mathrm{1}};\,\Varid{t}_{\mathrm{2}}}} & = &  \semapp{C}{t_1} \vmseq \semess{E}{u}{v}{t_2} \\
% Dot-dot
\rulename{sdotdot} &  \semess{E}{u}{v}{\ensuremath{\Varid{t}_{\mathrm{1}}\mathinner{\ldotp\ldotp}\Varid{t}_{\mathrm{2}}}}
  & = & \vminject \conBrack{u=v} \vmseq \semess{E}{p_1}{q_1}{t_1} \vmseq \semess{E}{p_2}{q_2}{t_2} \vmseq \\
  & & & \vmunions \{\;  \vmconcat [ \vminject \conBrack{v=q_1+i, i \leq q_2 - q_1}  \vb i \leftarrow [0..n] ]  \vb n \in \N \; \} \\
  & & & \pomHIDE \{p_1,q_1,p_2,q_2\} \hfill  p_1,q_1,p_2,q_2 \not\in \freevars{t_1,t_2,u,v} \\[1mm]

% Application and type
\rulename{sapp} & \semess{E}{u}{v}{\ensuremath{\Varid{t}_{\mathrm{1}} [\!\Varid{t}_{\mathrm{2}}]}}
& = &  \vminject \conBrack{u=v} \vmseq \semess{E}{f}{g}{t_1}
       \vmseq \semess{E}{p}{q}{t_2}
       \vmseq \semap{F}{\ensuremath{\Varid{g} [\!\Varid{q}]}}v  \\
    & & & \vmhide \{f,g,p,q\} \hfill f,g,p,q \not\in \freevars{t_1,t_2,u,v} \\[1mm]
\rulename{stype} &  \semess{E}{u}{v}{\ensuremath{\mathbin{:}\Varid{t}}} & = & \semess{E}{p}{q}{t} \vmseq \semap{F}{\ensuremath{\Varid{q} [\!\Varid{u}]}}v
  \vmhide \{p,q\} \hfill p,q \not\in \freevars{t,u,v} \\[1mm]

% IF
\multicolumn{2}{@{}l}{\rulename{sif} \; \hfill
  \semess{E}{u}{v}{\ensuremath{\textbf{if}(\Varid{t}_{\mathrm{0}})\{\Varid{t}_{\mathrm{1}}\}\!\;\textbf{else}\;\!\{\Varid{t}_{\mathrm{2}}\}}} }
& = & \begin{array}[t]{@{\;}l@{\,}l}
        (S0 \vmseq \semess{B}{u}{v}{t_1} \vmhide xs)
            \vmunion
          (\vmnot(S0 \vmhide xs) \vmseq \semess{B}{u}{v}{t_2}) \\
        \textbf{where} \; \xs = \semap{I}{\ensuremath{\Varid{t}_{\mathrm{0}}}}, \; S0 = \vmone( \semap{C}{\ensuremath{\Varid{t}_{\mathrm{0}}}}, xs )
      \end{array}
\\[5mm]

% FOR
\multicolumn{2}{@{}l}{\rulename{sfor} \; \hfill \semess{E}{u}{v}{\ensuremath{\textbf{for}(\Varid{t}_{\mathrm{0}})\,\textbf{do}\{\Varid{t}_{\mathrm{1}}\}}}}
& = &  \ldots \\[2mm]

% FUNCTIONS
  \rulename{sfun} &
   \semess{E}{h}{f}{\ensuremath{\textbf{fun}(\Varid{t}_{\Varid{a}})\langle\Varid{q}\rangle\langle\omega\rangle\{\Varid{t}_{\Varid{b}}\}}} \\
   & \multicolumn{3}{l}{\;\; = \vminject \begin{array}[t]{@{}l@{\;}l}
     \{  & \rho \in \Env \\
     \vb & \ensuremath{\mathbf{let}} \begin{array}[t]{l}
       \ensuremath{\Varid{dom}}:\Vm{\Env} \\
       \ensuremath{\Varid{dom}} = (\vminject\{ \rho \} \vmhide \semap{I}{\ensuremath{\Varid{t}_{\Varid{a}}}} \vmhide \{p,w\}) \vmseq \semess{E}{p}{w}{at}
       \end{array} \\
     , & \ensuremath{\mathit{fun}} : \Vm{\Val,\Val} \\
       & \ensuremath{\mathit{fun}} \in \vmpi \begin{array}[t]{@{}l@{\;}l}
               ( & \vmmap( \lambda \rho. \rho(p), \ensuremath{\Varid{dom}} ) \\
       , & \lambda val . \bigcap \begin{array}[t]{@{}l@{\;}l}
                            \{ & ((\{ \rho_0 \} \ENVhide \{j,z\} \ENVhide \semap{I}{t_b}) \cap \vmcollapse(\semess{E}{j}{z}{t_b}) \vmhide \semap{I}{t_a} \\
                            \vb & \rho_0 \in \vmcollapse(\ensuremath{\Varid{dom}}), \envlk{\rho_0}{p} = val \; \} \; )
                            \end{array}
               \end{array} \\
     , & \ensuremath{\mathbf{let}} \begin{array}[t]{l}
                   \ff, \hh :: \Vm{ \Val, \Val } \\
                   \ff = \vmmap( \; \lambda (\_,\rho). (\envlk{\rho}{p}, \envlk{\rho}{z}),\; \ensuremath{\mathit{fun}}\; ) \\
                   \hh = \vmmap( \; \lambda (\_,\rho). (\envlk{\rho}{w}, \envlk{\rho}{j}),\; \ensuremath{\mathit{fun}}\; ) \\
                \end{array} \\
     , & \envlk{\rho}{f} = \Fun(\ff) \\
     , & \envlk{\rho}{h} = \Fun(\hh) \; \} \\
  \end{array}
  }
   \\ \\
%%%%%%%%%%%%%%%%%%%%%%%%%%%%%
%       B, C, and F
%%%%%%%%%%%%%%%%%%%%%%%%%%%%%
& \semess{B}{u}{v}{\ensuremath{\Varid{t}}} & : & \VmENV \\
& \semess{B}{u}{v}{\ensuremath{\Varid{t}}} & = & \semess{E}{u}{v}{e} \pomHIDE \semapp{I}{\ensuremath{\Varid{t}}} \\[1mm]
& \semapp{C}{\ensuremath{\Varid{t}}} & : & \VmENV \\
& \semapp{C}{\ensuremath{\Varid{t}}} & = & \semess{E}{p}{q}{\ensuremath{\Varid{t}}} \pomHIDE \{p,q\} \hfill p,q \not\in \freevars{t}\\[1mm]
& \semap{F}{\ensuremath{\Varid{f} [\!\Varid{x}]}}r & : & \VmENV \\
& \semap{F}{\ensuremath{\Varid{f} [\!\Varid{x}]}}r  & = &
\vmunions \begin{array}[t]{@{}l}
           \,\{ \; \vmflatmap(\; \lambda (p,q). \conBrack{ f = \Fun(\ff), x=p, y=q }, \ff \; )\\
           \vb \Fun(\ff) \in \Val \; \}
           \end{array}
\end{array} \\
\\

\end{array}
$$
\caption{E-POM: denotational semantics of Essential Verse, using pomsets}
\label{fig:E-sls1}
\end{figure}


% --------------------------------------------
%                Functions
% --------------------------------------------

\begin{figure}
$$
\begin{array}{l}
\begin{array}{@{}lr@{\;\;}c@{\;\;}l}
                & \semess{E}{h}{f}{\ensuremath{\Varid{t}}} & : & \pomENV  \\
  \rulename{sfun} &
   \semess{E}{h}{f}{\ensuremath{\textbf{fun}(\Varid{t}_{\mathrm{0}})\langle\Varid{q}\rangle\langle\omega\rangle\{\Varid{t}_{\mathrm{1}}\}}}
    & = &  \semFUN{\semess{E}{p}{w}{\ensuremath{\Varid{t}_{\mathrm{0}}}}}{\semap{I}{t_0},q,\ensuremath{\omega}}{cd}{hf}{\ensuremath{\Varid{t}_{\mathrm{1}}}}
   \hspace{1cm} \hfill p,w,j,z \not\in \freevars{t_0,t_1,h,f,xs} \\
   & \multicolumn{3}{l}{\hspace{6mm}\textbf{where}} \\
   & \multicolumn{3}{l}{\hspace{1cm}
\begin{array}{@{}r@{\;\;}c@{\;\;}l}

  xs & = & \semap{I}{t_0} \\[1mm]
  \semFUN{\_}{xs,q,\ensuremath{\omega}}{cd}{hf}{t_1} & : & \pomENV \rightarrow \pomENV \\
% &  \semFUN{\pomempty}{xs}{cd}{hf}{t_1} & = & \pomunit{\conBrack{h=f=\Fun(\pomempty)}} \\[1mm]
% &  \semFUN{S \pomunion T}{xs}{cd}{hf}{t_1 } & = & \semFUN{S}{xs}{cd}{hf}{t_1} \pomunion \semFUN{T}{xs}{cd}{hf}{t_1} \\[1mm]

  \semFUN{S \pomunion T}{xs,q,\ensuremath{\omega}}{cd}{hf}{t_1}
                   & = & \pomunit{\conBrack{h=h_1 \fununion h_2, f=f_1 \fununion f_2}}
                   \hspace{6mm} \hfill h_1,f_1,h_2,f_2 \not\in \freevars{t_1,h,f,xs} \\
              &&  \pomseq \semFUN{S}{xs,q,\ensuremath{\omega}}{cd}{h_1f_1}{t_1} \\
              &&  \pomseq \semFUN{T}{xs,q,\ensuremath{\omega}}{cd}{h_2f_2}{t_1} \\
              && \pomHIDE \{h_1,f_1,h_2,f_2\} \\[1mm]

  \semFUN{S \pomcat T}{xs,q,\ensuremath{\omega}}{cd}{hf}{t_1}
                   & = & \pomunit{\conBrack{h=h_1 \funcat h_2, f=f_1 \funcat f_2}}
                   \hspace{3mm} \hfill h_1,f_1,h_2,f_2 \not\in \freevars{t_1,h,f,xs} \\
              &&  \pomseq \semFUN{S}{xs,q,\ensuremath{\omega}}{cd}{h_1f_1}{t_1} \\
              && \pomseq \semFUN{T}{xs,q,\ensuremath{\omega}}{cd}{h_2f_2}{t_1} \\
              && \pomHIDE \{h_1,f_1,h_2,f_2\} \\[1mm]

% --------- FUNCTIONS ATTEMPT 1 ---------------
 \text{Attempt OUROLD} \\
  \semFUN{\pomunit{\EV_0}}{xs}{cd}{hf}{t_1}
  &  = & \pomunit{
  \begin{array}[t]{@{}l@{\;}ll}
     \{ & \rho \in \Env \\
     \vb & \exists \hh, \EV_1. \\
     & \ensuremath{\mathbf{let}}\; \EV' = \EV_0 \cap (\rho \pomHIDE \{c,d,u,v,xs\}) & \rulename{(f1)} \\
     , & \Fun{\pomunit{\hh}} = \rho(h)  & \rulename{(f2)} \\
     , & \ensuremath{\textbf{if}\;\omega\mathrel{=}\langle\textbf{succeeds}\rangle}  & \rulename{(f3)} \\
      & \hspace{2mm} \ensuremath{\textbf{then}} \; \dom(\hh) = \envlk{\EV'}{d} \\
      & \hspace{2mm} \ensuremath{\textbf{else}} \; \dom(\hh) \subseteq \envlk{\EV'}{d}  \\
     , & \pomunit{\EV_1} = \pomunit{\EV' \cap \conBrack{ \fpr{d}{u} \in hh }}
                 \pomseq \semess{E}{u}{v}{\ensuremath{\Varid{t}_{\mathrm{1}}}} & \rulename{(f4)} \\
     , & \pomempty = \pomunit{\EV' \cap \conBrack{ d \not\in \dom(hh) }}  \pomseq \semess{E}{u}{v}{\ensuremath{\Varid{t}_{\mathrm{1}}}} & \rulename{(f5)} \\
     , & \ensuremath{\mathbf{let}}\; \ff = \{ \fpr{\envlk{\rho_1}{c}}{\envlk{\rho_1}{v}} \vb \rho_1 \in \EV_1 \} & \rulename{(f6)}\\
     , & \dom(\ff) = \{ \envlk{\rho'}{c} \vb \rho' \in \EV', \envlk{\rho'}{d} \in \dom(\hh) \} & \rulename{(f7)}\\
     , & isFunction(\ff) & \rulename{(f8)} \\
     , & \Fun{\pomunit{\ff}} = \rho(f)  \} & \rulename{(f9)}

  \end{array} \\
  }
  \\

% % --------- FUNCTIONS ATTEMPT 2 ---------------
%  \semFUN{\pomunit{\EV_0}}{xs}{cd}{hf}{t_1}
%  &  = & \pomunit{ \pomunions
%  \begin{array}[t]{@@{}l@@{\;}ll}
%    \{ & \conBrack{ h = \Fun{\pomunit{\hh}}, f = \Fun{\pomunit{\ff}}} \\
%    \vb & \exists \EV_1. \\
%    & \hh \in \pomFUN \\
%     , & |let|\; \EV' = \EV_0 \cap (\rho \pomHIDE \{c,d,u,v,xs\}) & \rulename{(f1)} \\
%     , & \Fun{\pomunit{\hh}} = \rho(h)  & \rulename{(f2)} \\
%     , & |if fx = eff succeeds|  & \rulename{(f3)} \\
%      & \hspace{2mm} |then| \; \dom(\hh) = \envlk{\EV'}{d} \\
%      & \hspace{2mm} |else| \; \dom(\hh) \subseteq \envlk{\EV'}{d}  \\
%     , & \pomunit{\EV_1} = \pomunit{\EV' \cap \conBrack{ \fpr{d}{u} \in hh }}
%                 \pomseq \semess{E}{u}{v}{|t_1|} & \rulename{(f4)} \\
%     , & \pomempty = \pomunit{\EV' \cap \conBrack{ d \not\in \dom(hh) }}  \pomseq \semess{E}{u}{v}{|t_1|} & \rulename{(f5)} \\
%     , & |let|\; \ff = \{ \fpr{\envlk{\rho_1}{c}}{\envlk{\rho_1}{v}} \vb \rho_1 \in \EV_1 \} & \rulename{(f6)}\\
%     , & \dom(\ff) = \{ \envlk{\rho'}{c} \vb \rho' \in \EV', \envlk{\rho'}{d} \in \dom(\hh) \} & \rulename{(f7)}\\
%     , & isFunction(\ff) & \rulename{(f8)} \\
%     , & \Fun{\pomunit{\ff}} = \rho(f)  \} & \rulename{(f9)}
% 
%  \end{array} \\
%  }
% \\
% --------- FUNCTIONS ATTEMPT 3 ---------------
\text{Attempt OURNEW} \\
  \semFUN{\pomunit{\EV_0}}{xs}{cd}{hf}{t_1}
  &  = & \pomunit{
  \begin{array}[t]{@{}l@{\;}ll}
     \{ & \rho \in \Env \\
     \vb & \ensuremath{\mathbf{let}}\; \EV_0' = \EV_0 \cap (\rho \pomHIDE \{p,q,j,z,xs\}) & \rulename{(f1)} \\
     , & \exists \EV_1. \pomunit{\EV_1} = \pomunit{\EV_0'} \pomseq \semess{E}{j}{z}{\ensuremath{\Varid{t}_{\mathrm{1}}}} & \rulename{(f2)} \\
     , & \ensuremath{\omega\mathrel{=}\langle\textbf{decides}\rangle} \; \vee \; \EV_0'(p) = \EV_1(p)  & \rulename{(f3)} \\
     , & \rho(f) = \Fun{ \pomunit{ \mkfunction \{ \fpr{ \rho_1(p)}{\rho_1(z)} \vb \rho_1 \in \EV_1 \} } }  & \rulename{(f4)} \\
     , & \rho(h) = \Fun{ \pomunit{ \mkfunction \{ \fpr{ \rho_1(w)}{\rho_1(j)} \vb \rho_1 \in \EV_1 \} } }  & \rulename{(f5)} \\
       \}

  \end{array} \\ }


\end{array}} \\
\end{array}
\\ \\
% \text{Note 1: for |eff succeeds|, \rulename{f5} is vacuously true } \\
% \text{Note 2: for |eff succeeds|, \rulename{f7} is equivalent to $\dom(\ff) = \envlk{\EV'}{c}$}

\end{array}
$$
\caption{E-POM: denotational semantics of Essential Verse: functions}
\label{fig:E-sls-funs}
\end{figure}



%%%%%%%%%%%%%%%%%%%%%%%%%%%%%%%%%%%%%
%     Commented out for now
%%%%%%%%%%%%%%%%%%%%%%%%%%%%%%%%%%%%%
\begin{comment}

% ------------------------------------------------------------------------------------
%
%           The SL monad
%
% ------------------------------------------------------------------------------------

\begin{figure}
$$
\begin{array}{l}
  \multicolumn{1}{c}{\fbox{\parbox{8cm}{
   \begin{center}
      ${\bf type} \; \SL{A} = \powersetnp{\seq{A}}$ \\
      \text{The set may be infinite, but the list is always finite}
   \end{center}
 }}} \\ \\

\textbf{SL is an applicative functor} \\[1mm]
\begin{array}{@{}r@{\;\;}c@{\;\;}l}
     \pureSL & : & \ENV \rightarrow \SL{\ENV}  \\
     \pureSL( \EV ) & = & \{ [ \EV ] \} \\[1mm]

\liftATwo & : & (A \rightarrow B \rightarrow C) \rightarrow \SL{A} \rightarrow \SL{B} \rightarrow\SL{C} \\
\liftATwo(OP,S_1,S_2) & = &
  \{ \; \concat(ss)
     \begin{array}[t]{@{}l@{\;}l}
     \vb & s_1 {:} \seq{A} \in S_1 \\
     , & \letvar t {:} \seq{\SL{C}} = [\; \mapSL(OP(\EV,\_), S_2) \vb \EV \ensuremath{\leftarrow } s_1 \; ] \\
     , & \letvar r {:} \SL{\seq{C}} = \pickE(t) \\
     , & ss {:} \seq{\seq{C}} \in r \; \}
     \end{array} \\
     \\[-2mm]

 (\seqSL) & : & \SL{\ENV} \rightarrow \SL{\ENV} \rightarrow \SL{\ENV} \\
 E_1 \seqSL E_2 & = & \liftATwo(\, \cap, \, E_1, \, E_2 \, ) \\[1mm]
 (\seqSLOne) & : & \ENV \rightarrow \SL{\ENV} \rightarrow \SL{\ENV} \\
 \EV_1 \seqSLOne E_2 & = & \pureSL(\EV_1) \seqSL E_2 \\
              & = & \mapSL(\;  \lambda \EV_2.\, \EV_1 \cap \EV_2, E_2 \;) \\[1mm]
  \mapSL & : & (A \rightarrow B) \rightarrow \SL{A} \rightarrow \SL{B} \\
  \mapSL(OP, E) & = & \{ \; [ OP( \EV) \vb \EV \ensuremath{\leftarrow } s ] \vb s \in E \; \} \\[1mm]
(\HIDESkw) & : &  \SL{\ENV} \rightarrow \powerset{\Ident} \rightarrow \SL{\ENV} \\
\HIDES{s}{xs} & = & \mapSL( \lambda \EV.\, \HIDE{\EV}{xs}, \; s) \\
\end{array}
\\
\\
\text{ \textbf{SL is also a monad (I think)} and bind is easier to think about than liftA2!} \\[1mm]
\begin{array}{@{}r@{\;\;}c@{\;\;}l}
(\bindSL) & : & \SL{A} \rightarrow (A \rightarrow \SL{B}) \rightarrow \SL{B} \\
S \bindSL K & = & \{ \; \concat(ss)
\begin{array}[t]{@{}l@{\;}l}
   \vb & s{:}[A] \in S \\
   , & \letvar t {:} [\SL{C}] = [\, K(\EV) \vb \EV \ensuremath{\leftarrow } s\, ] \\
   , & \letvar r {:} \SL{[C]} = \pickE(t) \\
   , & ss{:}[[C]] \in r \; \}
\end{array} \\
& = & \{ \; \concat(ss) \vb s \in S, \; ss \in \pickE[\, K(\EV) \vb \EV \ensuremath{\leftarrow } s\, ] \; \} \hspace{3mm} \text{\emph{more compactly}}
  \\[1mm]
\liftATwo(OP,S_1, S_2) & = & \begin{array}[t]{@{}l}
  S_1 \bindSL (\lambda \EV_1. \\
  S_2 \bindSL (\lambda \EV_2. \\
  \pureSL( OP(\EV_1,\EV_2) )\,))
  \end{array}
\end{array} \\
\\
\parbox{8cm}{NB: this naive version of $\liftATwo$ is wrong:
  $$\liftATwo(OP,S_1,S_2)
  = \{ [ OP(\EV_1,\EV_2) \vb \EV_1 \ensuremath{\leftarrow } s_1, \EV_2 \ensuremath{\leftarrow } s_2 ] \vb s_1 \in S_1, s_2 \in S_2 \}$$
...but there is no naive $\bindSL$.   True?}

\end{array}
$$
\caption{The $\SL{A}$ applicative functor}
\end{figure}

% ------------------------------------------------------------------------------------
%
%           Denotational semantics of Essential Verse in SLS style
%
% ------------------------------------------------------------------------------------

\begin{figure}
$$
\begin{array}{l}
\begin{array}{@{}lr@{\;\;}c@{\;\;}l}
    & \semess{E}{u}{v}{\ensuremath{\Varid{t}}} & : & \SL{\Env} \\

% Wildcards and atoms
\rulename{swild}  &  \semess{E}{u}{v}{\ensuremath{\anonymous }} & = & \pureSL \conBrack{u =v} \\
\rulename{svar}   &  \semess{E}{u}{v}{\ensuremath{\Varid{x}}} & = & \pureSL \conBrack{u=v=x} \\
\rulename{sconst} &  \semess{E}{u}{v}{\ensuremath{\Varid{k}}} & = & \pureSL \conBrack{u=v=k} \\
\rulename{sprim}  &  \semess{E}{u}{v}{\ensuremath{\Varid{o}}} & = & \pureSL \conBrack{u=v=o} \\

% Capture input and output
\rulename{sbind} &  \semess{E}{u}{v}{\ensuremath{\Varid{x}\mathrel{\coloneqq}\Varid{t}}} & = &  \conBrack{x=v} \seqSLOne \semess{E}{u}{v}{t} \\
\rulename{ssquig} &  \semess{E}{u}{v}{\ensuremath{\Varid{x}\!:\shortrightarrow\!\Varid{t}}} & = &
   \conBrack{x=u} \seqSLOne \semess{E}{u}{v}{t} \\

% Tuples
\multicolumn{2}{@{}l}{\rulename{stup} \; \hfill \semess{E}{u}{v}{\ensuremath{\textbf{array}\{\Varid{t}_{\mathrm{1}},\mycdots,\Varid{t}_{\Varid{n}}\}}}}
& = & \begin{array}[t]{@{}l}
         \conBrack{ u=\ensuremath{\langle\Varid{u}_{\mathrm{1}},\mycdots,\Varid{u}_{\Varid{n}}\rangle},
                    v=\ensuremath{\langle\Varid{v}_{\mathrm{1}},\mycdots,\Varid{v}_{\Varid{n}}\rangle}}
         \seqSLOne \semess{E}{u_1}{v_1}{t_1}
         \seqSL \mydots
         \seqSL \semess{E}{u_n}{v_n}{t_n} \\
         \HIDESkw { \{u_1,\mydots u_n,v_1 \mydots v_n\} }
         \hspace{5mm} u_i,v_i\not\in \freevars{t_1 \mydots t_n, u,v}
      \end{array} \\

% Choice and failure
\rulename{sfail} &  \semess{E}{u}{v}{\ensuremath{\textbf{fail}}} & = & \{ [ ] \} \\
\rulename{schoice} &  \semess{E}{u}{v}{\ensuremath{\Varid{t}_{\mathrm{1}}\hspace{0.19em}\mathop{\rule[-0.1em]{0.15em}{0.75em}}\hspace{0.2em}\Varid{t}_{\mathrm{2}}}}
  & = & \{ \; s_1 \append s_2 \vb s_1 \in \semess{B}{u}{v}{t_1}, s_2 \in \semess{B}{u}{v}{t_2} \; \} \\

% Unification
\rulename{sunify} &  \semess{E}{u}{v}{\ensuremath{\Varid{t}_{\mathrm{1}}\hspace{-0.14em}=\hspace{-0.14em}\Varid{t}_{\mathrm{2}}}} & = &   \semess{E}{u}{v} {t_1} \seqSL \semess{E}{u}{v}{t_2} \\

% Sequencing
\rulename{sseq} &  \semess{E}{u}{v}{\ensuremath{\Varid{t}_{\mathrm{1}};\,\Varid{t}_{\mathrm{2}}}} & = &  \semapp{C}{t_1} \seqSL \semess{E}{u}{v}{t_2} \\
% Dot-dot
\rulename{sdotdot} &  \semess{E}{u}{v}{\ensuremath{\Varid{a}_{\mathrm{1}}\mathinner{\ldotp\ldotp}\Varid{a}_{\mathrm{2}}}}
& = & \conBrack{u=v} \seqSLOne
   \{  \; [\conBrack{v=a_1+i, i \leq a_2 - a_1}  \vb i \leftarrow [0..n] ]  \vb n \in \N \} \\

% Application and type
\rulename{sapp} & \semess{E}{u}{v}{\ensuremath{\Varid{a}_{\mathrm{1}} [\!\Varid{a}_{\mathrm{2}}]}}
 & = &  \conBrack{u=v} \seqSLOne \semap{F}{\ensuremath{\Varid{a}_{\mathrm{1}} [\!\Varid{a}_{\mathrm{2}}]}}v \\
\rulename{stype} &  \semess{E}{u}{v}{\ensuremath{\mathbin{:}\Varid{a}}} & = & \semap{F}{\ensuremath{\Varid{a} [\!\Varid{u}]}}v  \\
\\

% IF
\multicolumn{2}{@{}l}{\rulename{sif} \; \hfill
  \semess{E}{u}{v}{\ensuremath{\textbf{if}(\Varid{t}_{\mathrm{0}})\{\Varid{t}_{\mathrm{1}}\}\!\;\textbf{else}\;\!\{\Varid{t}_{\mathrm{2}}\}}} }
& = &   \bigcup\limits_{s_0 \in \semap{C}{\ensuremath{\Varid{t}_{\mathrm{0}}}}}\left\{
 \begin{array}{@{\;}l@{\,}l}
(\HIDES{\EV_1 \seqSLOne \semess{B}{u}{v}{t_2})}{xs} \\
\cup \\
(\HIDE{\conCompl(\EV_1)}{xs}) \seqSLOne \semess{B}{u}{v}{t_3} \\
 \textbf{where} \; \xs = \semap{I}{\ensuremath{\Varid{t}_{\mathrm{0}}}}, \; \EV_1 = \firstE(xs,s_0)
 \end{array}
 \right\}
\\ \\
\multicolumn{2}{@{}l}{\rulename{sfor} \; \hfill \semess{E}{u}{v}{\ensuremath{\textbf{for}(\Varid{t}_{\mathrm{1}})\,\textbf{do}\{\Varid{t}_{\mathrm{2}}\}}}}
& = &   \bigcup\limits_{s_0 \in \semap{C}{\ensuremath{\Varid{t}_{\mathrm{0}}}}} \left\{
 \begin{array}{@{\;}l@{\,}l}
   \conBrack{ u=(u_1 \oplus \mydots \oplus u_n), v=(v_1 \oplus \mydots \oplus v_n)}
   \seqSLOne D_1 \seqSL \mydots \seqSL D_n \\
\hspace{5mm} \HIDESkw \{u_1, \mydots, u_n, v_1, \mydots, v_n \} \\
\textbf{where} \\
\;  \begin{array}[t]{l}
  \{u_1, \mydots, u_n, v_1, \mydots, v_n \} \not\in \freevars{u,v,t_0,t_1} \\
  {} [\EV_1,\mydots,\EV_n] = s_0 \\
  S1 = \semess{B}{p}{q}{t_1} \\
  C_i = \begin{array}[t]{@{}l}
      \HIDES{(\EV_i \seqSLOne \conBrack{u_i=\ensuremath{\langle\Varid{p}\rangle}, v_i=\ensuremath{\langle\Varid{q}\rangle}} \seqSLOne S1)}{\{p,q\}} \\
      \cup \\
      \conCompl( \EV_i ) \seqSLOne \pureSL \conBrack{u_i=\ensuremath{\langle\rangle}, v_i=\ensuremath{\langle\rangle}}
      \end{array} \\
  D_i = \HIDES{C_i}{xs} \\ 
  \xs = \semap{I}{\ensuremath{\Varid{t}_{\mathrm{0}}}}
 \end{array}
 \end{array}
 \right\}
\end{array}
\\

\\
%%%%%%%%%%%%%%%%%%%%%%%%%%%%%
%       B, C, and F
%%%%%%%%%%%%%%%%%%%%%%%%%%%%%
\begin{array}{@{}lr@{\;\;}c@{\;\;}l}
& \semess{B}{u}{v}{\ensuremath{\Varid{t}}} & : & \SL{ \ENV } \\
& \semess{B}{u}{v}{\ensuremath{\Varid{t}}} & = & \HIDES{\semess{E}{u}{v}{e}}{\semapp{I}{\ensuremath{\Varid{t}}}} \\[1mm]
& \semapp{C}{\ensuremath{\Varid{t}}} & : & \SL{ \ENV } \\
& \semapp{C}{\ensuremath{\Varid{t}}} & = & \HIDES{\semess{E}{p}{q}{\ensuremath{\Varid{t}}}}{\{p,q\}} \hspace{5mm} p,q \not\in \freevars{t}\\[1mm]
& \semap{F}{\ensuremath{\Varid{a}_{\mathrm{1}} [\!\Varid{a}_{\mathrm{2}}]}}r & : & \SL{\ENV} \\
& \semap{F}{\ensuremath{\Varid{a}_{\mathrm{1}} [\!\Varid{a}_{\mathrm{2}}]}}r  & = &
    \{ \ [ \; \conBrack{\conFun{a_2}{r}{a_1}{i}{n}} \vb i \ensuremath{\leftarrow } [0..n] \; ] \vb n \in \N
      \; \}
 \\[1.5ex]

& \firstE & : & \powerset{\Ident} \rightarrow [\ENV] \rightarrow \ENV \\
& \firstE(xs,\; []) & = & \{\} \\
& \firstE(xs,\; \EV : s)  & = & \EV \cup (\firstE(xs, s) \cap \conCompl(\HIDE{\EV}{xs})) \\
\end{array}
\end{array}
$$
\caption{E-SLS: denotational semantics of Essential Verse, in SLS style}
\label{fig:E-sls1}
\end{figure}

%%%%%%%%%%%%%%%%%%%%%%%%%%%%%%%%%%%%%
%     End of commented out for now
%%%%%%%%%%%%%%%%%%%%%%%%%%%%%%%%%%%%%
\end{comment}
